\subsection{Validation under YCSB-Driven Workloads}
\label{sec:eval-ycsb}

This section introduces segment-size skew as a robustness
dimension absent from the synthetic experiments above.
We use a YCSB Zipfian distribution ($\alpha{=}0.99$) to
generate the base keyspace and evaluate the following fault families:
\emph{Bimodal} (bounded-blast-radius concentrated faults),
\emph{Drift-shaped} (lag concentrated on hot segments),
and \emph{Uniform Diffuse} (broad low-rate anomalies).
Each family is instantiated at three data-scale tiers---Light
(20M records, $P{=}100{,}000$), Medium (200M), and Heavy
(2B)---which induce progressively denser segment fill states
under the same Zipfian workload.
A fourth family, \emph{Weighted Diffuse} (per-record loss
under segment-size skew), is further detailed in
Supplement~\ref{sec:appendix-ycsb-weighted} along with
parameterization details for all families
(Supplement Tables~\ref{tab:c6-trace-density}
and~\ref{tab:appendix-c6-faults}).

\paragraph{Gating}
The upper panel of Table~\ref{tab:c6-trace-triage} reports
SCOUT's $K_{\text{eff}}/L$, as \S\ref{sec:eval-comparison}
shows that neither baseline provides a reliable gating signal.
All four fault families form a clear gradient: Bimodal at
approximately $0.003$ (highly concentrated), Drift-shaped at
approximately $0.08$, Weighted Diffuse at approximately
$0.12$, and Uniform Diffuse at approximately $0.93$ (highly
diffuse). The wide separation between the first three
families ($K_{\text{eff}}/L \le 0.12$) and Uniform Diffuse
(${\ge}0.93$) leaves a broad margin for placing
$\tau_{\text{gating}}$: the first three families are admitted to the triage path while Uniform Diffuse is diverted to fallback.
$K_{\text{eff}}/L$ is nearly invariant across the three
data-scale tiers, indicating that the gating metric is
governed by segment-level fault shape rather than data
density. All families yield sufficient occupancy for decisive
gating across all three tiers, with no \texttt{unknown}
outcomes; full distributions are shown in
Supplement Figure~\ref{fig:c6-gating}.

\paragraph{Triage}
The lower panel of Table~\ref{tab:c6-trace-triage} reports the
minimum candidate fraction achieving
$\texttt{recall\_big}\ge 0.95$ under a 2\% cap. Bimodal faults
yield a stable candidate fraction of $3{\times}10^{-5}$;
Drift-shaped faults yield approximately
$7.3{\times}10^{-4}$. Across the three tiers, candidate
fractions change little: with $P$ and the segment-level fault
pattern held fixed, increasing the record count mainly
densifies segment fill rather than changing the sketch-level
concentration pattern. Uniform Diffuse does not appear in the
triage panel because gating diverts it to fallback.

\paragraph{Sizing}
Drift-shaped faults are the hardest sizing case because they
combine segment-size skew with concentration on hot segments.
With the Bessel correction applied, the de-drifted estimator
achieves a median relative error of 0.8\%--1.4\% across
the three tiers (p90 $\le 2.8\%$), matching the accuracy
observed under the uniform settings of
\S\ref{sec:eval-dedrift}; raw-control results are reported in
Supplement Table~\ref{tab:c6-sizing}.

\begin{table}[!thbp]
\caption{YCSB-driven validation summary across data-scale tiers.
Upper: gating statistic; lower: triage candidate fraction.}
\label{tab:c6-trace-triage}
\centering
\small
\setlength{\tabcolsep}{3pt}
\begin{tabular}{lccc}
\toprule
Family & Light & Medium & Heavy \\
\midrule
\multicolumn{4}{c}{SCOUT Gating: median $K_{\text{eff}}/L$} \\
\midrule
Bimodal          & 0.003 & 0.003 & 0.003 \\
Drift-shaped     & 0.082 & 0.082 & 0.081 \\
Weighted Diffuse & 0.122 & 0.118 & 0.120 \\
Uniform Diffuse  & 0.931 & 0.938 & 0.939 \\
\midrule
\multicolumn{4}{c}{SCOUT Triage: min $\texttt{cand\_frac}$
  at $\texttt{recall\_big}{\ge}0.95$} \\
\midrule
Bimodal      & $3{\times}10^{-5}$ & $3{\times}10^{-5}$ & $3{\times}10^{-5}$ \\
Drift-shaped & $7.3{\times}10^{-4}$ & $7.4{\times}10^{-4}$ & $7.3{\times}10^{-4}$ \\
\bottomrule
\end{tabular}
\end{table}

Sections~\ref{sec:eval-spatial}--\ref{sec:eval-ycsb} evaluate SCOUT's triage, sizing, and gating behavior across controlled synthetic workloads and YCSB-driven skewed workloads. Because
these results assume uniform hashing at the probe interface, we additionally verify in Supplement~\ref{sec:appendix-hash} that the triage operating points remain stable across three hash
implementations, two key structures, and $k\in\{1,2,3\}$.
\phantomsection\label{sec:eval-summary}
